% GENERATED FILE. Edit canonical lessons, then run npm run generate:books.
% canonical-source-hash: fnv1a64:856c6f564ba90029
% canonical-lessons: LA-C511-crudelitas, LA-C511-saevitia, LA-C511-laedo, LA-C511-mitigo, LA-C511-mansuetus

\chapter{Cruelty and gentleness}
\label{ch:la-cruelty-and-gentleness}
\hlchaptermodality{511}

\begin{chapteropening}
\textbf{By the end of this chapter:} \emph{I can talk about cruelty and gentleness.}

Complete the last lesson of chapter 511: I can talk about cruelty and gentleness.
\end{chapteropening}

\section[crūdēlitās, crūdēlitātis]{crūdēlitās, crūdēlitātis — cruelty}

\label{lesson:LA-C511-crudelitas}



\begin{quote}
Before the new one: say the Latin for out of one's mind, then the Latin for to be delirious.
\end{quote}

\subsection*{You'll want to know: crūdēlitās, crūdēlitātis}
\textbf{crūdēlitās, crūdēlitātis} — \textquotedblleft{}cruelty\textquotedblright{}.

\begin{quote}
Crūdēlitās dominī servōs terret.
\end{quote}

\begin{quote}
\emph{The master's cruelty frightens the slaves.}
\end{quote}

Feminine, third declension, from \emph{crūdēlis}. English \emph{cruelty}.

\begin{cousinweb}
\textbf{crūdēlitās} is built on \emph{crūdēlis} \textquotedblleft{}cruel\textquotedblright{}, itself related to \emph{crūdus} \textquotedblleft{}raw\textquotedblright{} (Lewis \& Short). English \emph{cruelty} comes from it.
\end{cousinweb}

\begin{grammarlens}[title={What you've built}]
One more word for telling, explaining and describing at B1.
\end{grammarlens}

\subsection*{Your turn}
\begin{itemize}
\raggedright
  \item \emph{Say it:} \emph{crūdēlitās, crūdēlitātis}
  \item \emph{Say it:} \emph{crūdēlitās, crūdēlitātis}, once more
  \item \emph{Say it:} a sentence of your own with \emph{crūdēlitās, crūdēlitātis}
  \item \emph{Recall:} say the Latin for out of one's mind, then the Latin for to be delirious, then say \emph{crūdēlitās, crūdēlitātis} again
\end{itemize}

\begin{tcolorbox}[breakable,colback=teal!4,colframe=teal!35!black,title={Before you move on}]
What does crūdēlitās, crūdēlitātis mean? (\textquotedblleft{}Cruelty\textquotedblright{}.) Say it once more.
\end{tcolorbox}
\section[saevitia, saevitiae]{saevitia, saevitiae — fury, savagery}

\label{lesson:LA-C511-saevitia}



\begin{quote}
Before the new one: say the Latin for to be delirious, then the Latin for cruelty.
\end{quote}

\subsection*{You'll want to know: saevitia, saevitiae}
\textbf{saevitia, saevitiae} — \textquotedblleft{}fury, savagery\textquotedblright{}.

\begin{quote}
Saevitia tempestātis nautās terret.
\end{quote}

\begin{quote}
\emph{The fury of the storm frightens the sailors.}
\end{quote}

Feminine, first declension, from \emph{saevus} — fierce. Of beasts, storms and harsh masters.

\begin{cousinweb}
\textbf{saevitia} is built on \emph{saevus} \textquotedblleft{}fierce, savage\textquotedblright{} (Lewis \& Short; EtymDB); the origin of \emph{saevus} is not settled.
\end{cousinweb}

\begin{grammarlens}[title={What you've built}]
One more word for telling, explaining and describing at B1.
\end{grammarlens}

\subsection*{Your turn}
\begin{itemize}
\raggedright
  \item \emph{Say it:} \emph{saevitia, saevitiae}
  \item \emph{Say it:} \emph{saevitia, saevitiae}, once more
  \item \emph{Say it:} a sentence of your own with \emph{saevitia, saevitiae}
  \item \emph{Recall:} say the Latin for to be delirious, then the Latin for cruelty, then say \emph{saevitia, saevitiae} again
\end{itemize}

\begin{tcolorbox}[breakable,colback=teal!4,colframe=teal!35!black,title={Before you move on}]
What does saevitia, saevitiae mean? (\textquotedblleft{}Fury, savagery\textquotedblright{}.) Say it once more.
\end{tcolorbox}
\section[laedō, laedere]{laedō, laedere — to hurt, harm; to offend}

\label{lesson:LA-C511-laedo}



\begin{quote}
Before the new one: say the Latin for cruelty, then the Latin for fury.
\end{quote}

\subsection*{You'll want to know: laedō, laedere}
\textbf{laedō, laedere} — \textquotedblleft{}to hurt, harm; to offend\textquotedblright{}.

\begin{quote}
Verba tua amīcum laedunt.
\end{quote}

\begin{quote}
\emph{Your words hurt your friend.}
\end{quote}

\emph{laedō, laedere, laesī, laesum}. Of the body and of feelings. English \emph{lesion}.

\begin{cousinweb}
\textbf{laedō} is a native Latin verb; Lewis \& Short offers a guess at its root, but the origin is not settled. English \emph{lesion} comes from its participle \emph{laesus}, and \emph{collide} from the compound \emph{collīdō}.
\end{cousinweb}

\begin{grammarlens}[title={What you've built}]
One more word for telling, explaining and describing at B1.
\end{grammarlens}

\subsection*{Your turn}
\begin{itemize}
\raggedright
  \item \emph{Say it:} \emph{laedō, laedere}
  \item \emph{Say it:} \emph{laedō, laedere}, once more
  \item \emph{Say it:} a sentence of your own with \emph{laedō, laedere}
  \item \emph{Recall:} say the Latin for cruelty, then the Latin for fury, then say \emph{laedō, laedere} again
\end{itemize}

\begin{tcolorbox}[breakable,colback=teal!4,colframe=teal!35!black,title={Before you move on}]
What does laedō, laedere mean? (\textquotedblleft{}To hurt, harm; to offend\textquotedblright{}.) Say it once more.
\end{tcolorbox}
\section[mītigō, mītigāre]{mītigō, mītigāre — to ease, soften, calm}

\label{lesson:LA-C511-mitigo}



\begin{quote}
Before the new one: say the Latin for fury, then the Latin for to hurt.
\end{quote}

\subsection*{You'll want to know: mītigō, mītigāre}
\textbf{mītigō, mītigāre} — \textquotedblleft{}to ease, soften, calm\textquotedblright{}.

\begin{quote}
Vīnum dolōrem mītigat.
\end{quote}

\begin{quote}
\emph{Wine eases pain.}
\end{quote}

\emph{mītigō, mītigāre, mītigāvī, mītigātum}: \emph{mītis} 'mild' + \emph{agō} 'to make'. English \emph{mitigate}.

\begin{cousinweb}
\textbf{mītigō} is built on \emph{mītis} \textquotedblleft{}mild\textquotedblright{} and \emph{agō} \textquotedblleft{}to drive, make\textquotedblright{} (Lewis \& Short; EtymDB). English \emph{mitigate} comes from it.
\end{cousinweb}

\begin{grammarlens}[title={What you've built}]
One more word for telling, explaining and describing at B1.
\end{grammarlens}

\subsection*{Your turn}
\begin{itemize}
\raggedright
  \item \emph{Say it:} \emph{mītigō, mītigāre}
  \item \emph{Say it:} \emph{mītigō, mītigāre}, once more
  \item \emph{Say it:} a sentence of your own with \emph{mītigō, mītigāre}
  \item \emph{Recall:} say the Latin for fury, then the Latin for to hurt, then say \emph{mītigō, mītigāre} again
\end{itemize}

\begin{tcolorbox}[breakable,colback=teal!4,colframe=teal!35!black,title={Before you move on}]
What does mītigō, mītigāre mean? (\textquotedblleft{}To ease, soften, calm\textquotedblright{}.) Say it once more.
\end{tcolorbox}
\section[mānsuētus, mānsuēta, mānsuētum]{mānsuētus, mānsuēta, mānsuētum — tame; gentle}

\label{lesson:LA-C511-mansuetus}



\begin{quote}
Before the new one: say the Latin for to hurt, then the Latin for to ease.
\end{quote}

\subsection*{You'll want to know: mānsuētus, mānsuēta, mānsuētum}
\textbf{mānsuētus, mānsuēta, mānsuētum} — \textquotedblleft{}tame; gentle\textquotedblright{}.

\begin{quote}
Leō in amphitheātrō mānsuētus est.
\end{quote}

\begin{quote}
\emph{The lion in the amphitheatre is tame.}
\end{quote}

The participle of \emph{mānsuēscō}: \emph{manus} 'hand' + \emph{suēscō} 'to get used to' — used to the hand.

\begin{cousinweb}
\textbf{mānsuētus} is the participle of \emph{mānsuēscō} \textquotedblleft{}to grow tame\textquotedblright{}, built on \emph{manus} \textquotedblleft{}hand\textquotedblright{} and \emph{suēscō} \textquotedblleft{}to become used to\textquotedblright{} (Lewis \& Short): used to the hand.
\end{cousinweb}

\begin{grammarlens}[title={What you've built}]
One more word for telling, explaining and describing at B1.
\end{grammarlens}

\subsection*{Your turn}
\begin{itemize}
\raggedright
  \item \emph{Say it:} \emph{mānsuētus, mānsuēta, mānsuētum}
  \item \emph{Say it:} \emph{mānsuētus, mānsuēta, mānsuētum}, once more
  \item \emph{Say it:} a sentence of your own with \emph{mānsuētus, mānsuēta, mānsuētum}
  \item \emph{Recall:} say the Latin for to hurt, then the Latin for to ease, then say \emph{mānsuētus, mānsuēta, mānsuētum} again
\end{itemize}

\begin{tcolorbox}[breakable,colback=teal!4,colframe=teal!35!black,title={Before you move on}]
What does mānsuētus, mānsuēta, mānsuētum mean? (\textquotedblleft{}Tame; gentle\textquotedblright{}.) Say it once more.
\end{tcolorbox}
